\section{Transporting the deck sign through the marking}
\label{sec:marked-sheet-transport}

This section keeps track of the sign lost by projectivizing the Clebsch
chart.  The normalization is intrinsic to the pullback.  Its comparison with
the marked icosahedral six-axis data and the Petersen coefficient module is relative to
the marking datum \(\mathfrak m\) fixed in the introduction.

\subsection{Normalization over the Clebsch chart}

Let \(\mathcal N\to\mathbf P(H)\) be the incidence normalization from
Section~\ref{sec:cover}, and put \(B=\mathbf P(V_E)\).  Choose the
anti-invariant generator in Theorem~\ref{thm:arithmetic-main}; its pullback
is the odd generator \(z\) below.  On the chart, the Stein-algebra theorem
and Hitchin's restriction identity give the equality of finite
\(\mathcal O_B\)-algebras
\[
 z^2=5\iota_t^*J_0=80\sigma_3^2
   =(z-4\sqrt5\,\sigma_3)(z+4\sqrt5\,\sigma_3).
\]
For a normal domain \(S\), with \(2\) invertible and \(0\ne a\in S\), the algebra
\(S[z]/(z^2-a^2)\) maps to \(S\times S\) by \(z\mapsto(a,-a)\).  Its image
consists of pairs congruent modulo \(a\); hence its normalization is
\(S\times S\), its conductor is \((a)\times(a)\), and it is already split
etale where \(a\) is a unit.  Applying this elementary pinching lemma with
\(a=4\sqrt5\,\sigma_3\) gives the normalization
\(B_+\amalg B_-\) and locates the entire failure of normality on the
Clebsch cubic.  The constant plane \(U_t\) supplies one component; deck
exchange supplies the other.  After base change to \(\mathbf C\), write
\(\Delta_t=\iota_t^*\Delta\) for the pullback of Hitchin's degree-ten
discriminant introduced in Section~\ref{sec:rational-incidence-model}; only
its zero locus is used.  The algebra above shows that on \(D(\sigma_3)\) the
pulled-back finite cover is already split etale, while its normalization
\(B_+\amalg B_-\) remains disjoint everywhere.  Hitchin
\cite[\S8.2]{HitchinBundles} identifies \(\{\Delta_t=0\}\) with the locus
where a cubic in the fixed Clebsch chart is also orthogonal to a degenerate
isotropic three-plane.  Hence the literal interpretation by two regular
six-axis configurations is valid precisely on \(D(\sigma_3\Delta_t)\).  On
\(D(\sigma_3)\cap\{\Delta_t=0\}\), the \(U_t\)-branch is regular and the
other compact incidence point is degenerate.
The magnitude \(4\sqrt5\) is forced by the product of the descent class
\(5\) and the chart factor \(16\); choosing its sign is precisely choosing
one normalized component.

The sign of \(z\) is the remaining freedom in the chosen anti-invariant
generator.  We normalize it by declaring the component supplied by \(U_t\)
to have \(z=4\sqrt5\,\sigma_3\).  This does not strengthen the integral
statement: comparison with the geometric incidence model still holds only
over an unspecified \(\mathbf Z[1/N]\).

\subsection{The configuration exchanger and global negation}

Choose the displayed normalized representatives
\(a_0,\ldots,a_5\) of the ordered projective axes in \(\mathfrak m\), and
write \(C=C_{\mathfrak m}\) and \(Z_C=Z_{\mathfrak m}\).  Their Gram
matrix has the form
\[
 G=(t+2)I+tC,
\qquad
C=\begin{pmatrix}
0&1&1&1&-1&-1\\
1&0&-1&-1&-1&-1\\
1&-1&0&1&1&-1\\
1&-1&1&0&-1&1\\
-1&-1&1&-1&0&-1\\
-1&-1&-1&1&-1&0
\end{pmatrix}.
\]
The six vertex axes of an icosahedron as equiangular lines, and their signed
Gram or Seidel-matrix description up to diagonal switching, are classical
\cite[\S3]{GodsilProblems}.  What is tracked below is the sign of the
specific marked triangle cubic under the configuration exchanger.
The frame operator \(\sum_i a_ia_i^{\mathsf T}\) commutes with the
irreducible three-dimensional \(A_5\)-action, so it is scalar.  Its trace is
\(6(t+2)\); division by the dimension gives
\(\sum_i a_ia_i^{\mathsf T}=2(t+2)I_3\).  Hence
\(G^2=2(t+2)G\).  Substituting \(G=(t+2)I+tC\) and using
\((t+2)^2=5t^2\) gives \(C^2=5I\).  Changing signs of axis
representatives switches \(C\) to \(DCD\), but every triangle product and
hence \(Z_C\) is unchanged.  The off-diagonal entries of \(C^2\) also give
\[
 \sum_{k\ne i,j}C_{ij}C_{jk}C_{ki}=0.
\]
Thus the directional derivative of \(Z_C\) along \(\mathbf1\) vanishes, so
\(Z_C(x+u\mathbf1)=Z_C(x)\) and the cubic descends to
\(\Q^6/\Q\mathbf1\).

Let \(a_i'\) denote the conjugate axes, obtained by \(t\mapsto1-t\), and
let \(R\) be the exchanger of Section~\ref{sec:arithmetic}.  Direct
substitution gives
\[
 Ra_i=t d_i a'_{p(i)},\qquad
 p=(0,1,5,4,2,3),\qquad
 (d_i)=(1,-1,1,1,1,1).
\]
Using \(t^2(1-t)=-t\), one obtains
\[
 d_id_jC_{p(i)p(j)}=-C_{ij}.
\]
Thus, over the incidence fibre at \([xyz]\), Galois conjugation transported
by \(R\) changes \(C\)
to \(-C\) and \(Z_C\) to \(-Z_C\).  Consequently, deck exchange leaves the unmarked
quadratic spectral algebra unchanged, \(\Q[C]=\Q[-C]\), even though it reverses
the chosen marked generator and the relative orientation data.  The lone
negative \(d_i\) is only a change of axis representatives and does not cause
the reversal; switching alone leaves the triangle cubic fixed.

No global marking of the varying configurations on \(B_-\) is needed or
asserted.  The fixed section \(U_t\) selects \(B_+\), but does not by itself
choose \(\mathfrak m\).  After \(\mathfrak m\) is fixed, we attach its
constant marked conference pair
\(\bigl(C_{\mathfrak m},Z_{\mathfrak m}\bigr)\) to \(B_+\) and attach the
deck-opposite pair to \(B_-\), meaning that the normalized
odd generator changes sign.  The calculation at \([xyz]\) proves that this
relative convention agrees with the actual Galois-conjugate icosahedral configuration
there; it is not used to infer a pointwise conference labeling elsewhere.
These relative pair labels are attached to the normalized components, not
to a pointwise regular-six-axis marking.  They therefore extend across both
\(\sigma_3=0\) and, after base change to \(\mathbf C\), \(\Delta_t=0\); the
literal interpretation by two regular icosahedral six-axis configurations is
valid only on \(D(\sigma_3\Delta_t)\).

\subsection{Transport to degree six}

Write Hitchin's chart as
\(\iota_t(y)=\sum_i y_iq_i^{(t)}\), with \(q_1=xyz\).  Since
\(R(xyz)=-xyz\), equivariance and irreducibility of the four-module give
\[
 Rq_i^{(t)}=-q_i^{(1-t)}
\]
after transporting the five plane-triple labels.  Normalize the lift on
\(B_+\) by \(q_1=xyz\); the deck-opposite pair carries the lift
\(-\iota_t\).  The displayed exchanger calculation shows that this
normalization agrees with Galois conjugation at \([xyz]\).  Thus the sign of
the lift is part of the marking datum, rather than a purported global marking
of the second family.

The primitive integral map
\[
 \beta(y)_{ij}=y_i+y_j
\]
lands in the Petersen \((-2)\)-eigenspace and satisfies
\[
 \sum_{i<j}(y_i+y_j)^2=3\sum_i y_i^2,
 \qquad
 y_i=\frac13\sum_{j\ne i}\beta(y)_{ij}.
\]
Multiplicity one makes this comparison unique after the five labels and the
normalization are fixed.  Reversing the chart lift reverses \(y\), the zonal harmonic
\(F_y\), and its cubic integral.  The two signs in the normalized incidence
cover therefore select the two signs of the exact Gaunt cubic relative to
\(\mathfrak m\).  This proves Proposition~\ref{thm:marked-sheet-comparison}.

The remaining marking actions and the boundary of this comparison are listed
in Appendix~\ref{app:marking-ambiguities}.
